\lecture{4}{Fri 22 Sep 2023 10:58}{Uniform Integrability of Martingale}

Motivation. Can every martingale be written as conditional expectations of a single r.v. with respect to a filtration? Answer: true for $L^p$ martingales, and more generally, for uniformly integrable martingales.

Recall: $L^p \implies \text{(Jensen)} L^1 \implies\text{(Chebyshev)} P \implies \text{weak}$.

\begin{theorem}
	Let $X_n$ be any set of r.v.s in $L^2$. Then $X_n \to X$ iff $X_n \to X$ in probability, and $\{X_n\} $ is uniformly integrable.
\end{theorem}

\begin{theorem}
	For a submartingale $\{M_n\} $ the following are equivalent:
	\begin{enumerate}
		\item It is uniformly integrable.
		\item It converges a.s. and in  $L^1$.
		\item It converges in $L^1$.
		\item (For martingales) There exists a r.v. $X \in (\Omega, \mathcal{F}, P)$ such that
			\[
				M_n = E(X | \mathcal{F}_n)
			.\] 
	\end{enumerate}
\end{theorem}
\begin{definition}
	A collection of r.v.s $\mathcal{L}$ is called uniformly integrable if given $s > 0, \mathcal{F}, M$ such that for all $X \in \mathcal{L}$, $E(|X|; |X| > M) < \varepsilon$.
\end{definition}

\begin{example}
	Take $\mathcal{L} = \{X\} , X \in L^1$. Chebyshev's inequality gives $P(|X| > \lambda) \le  \frac{1}{\lambda} E|X| \to  0$ as $\lambda \to \infty $.
	

	In general, for $f \in L^1(\mu)$, this means given $\varepsilon>0$, there exists $\lambda>0$ such that
	\[
		\int _A |f| d \mu < \varepsilon \impliedby \mu(A) < \lambda
	.\] 
	which is absolute continuity of measure.

	Set $A = \{|X| > M\} $.
\end{example}
\begin{example}
	(Durett 4.6.1)
	Suppos $X \in \mathcal{F}$ and $L^1$. 

	Take $\mathcal{L} = \{E(X | \mathcal{F}) : \mathcal{F_0} \subset \mathcal{F}\} $, then this collection is uniformly integrable.
\end{example}
\begin{proof}
	By using Jensen's Inequality we can have a uniform bound for $E\left(|E(X|\mathcal{F})| ; |E(X|\mathcal{F})| > M\right)$

	\[
		E\left(|E(X|\mathcal{F})| ; |E(X|\mathcal{F})| > M\right)
		\le E\left(E(|X||\mathcal{F}) ; |E(X|\mathcal{F})| > M\right)
		\le E\left(|X| ; |E(X|\mathcal{F})| > M\right)
	.\] 
	Next we provide a uniform bound on the event where expectation is taken. Using Chebyshev inequality (and tower property),
	\[
		P(|E(X|\mathcal{F})| > M)
		\le \frac{1}{M} E|X|
	.\] 
	Now as $M \to \infty $, for any $\mathcal{F}$, $P(|E(X|\mathcal{F})| > M) \to 0 $ and the convergence is uniform in $\mathcal{F}$. Hence $E\left(|X| ; |E(X|\mathcal{F})| > M\right) \to 0$ uniformly in $\mathcal{F}$, since $X$ is $L^1$, by DCT. This gives us uniform integrability.
\end{proof}


\begin{proposition}
Suppose $\mathcal{L}$ is uniformly integrable, then it is $L^1$ bounded.
\end{proposition}
\begin{proof}
	

Let $\varepsilon = 1$, then 
\[
	E|X| = E(|X|; |X| > M) + E(|X|; |X|<M)
.\] 
Then there exists a $M$ such that
\[
	E(|X|; |X| > M) \le 1; E(|X|; |X| < M) \le M
.\] 
Thus $E|X| \le  1 + M$ for all $X \in \mathcal{L}$.
\end{proof}

$L^1$ bounded does \textbf{not} imply uniform integrability.
\begin{example}
	Take $A_n = (0, \frac{1}{n})$ and $X_n = n 1_{A_n}$. Then $EX_n = 1$, but not converge to $0$ for any truncation $M$.
\end{example}

\begin{proposition}(Couple of sufficient conditions)
	\begin{enumerate}
		\item $1 < p < \infty $, $\mathcal{L} = L^p(P)$, $\|X\|_p \le  k < \infty $ for all $X \in L^p$, then $\mathcal{L}$ is uniformly integrable.
	\end{enumerate}
\end{proposition}
Simple application of Holder Inequality.

\begin{proposition}
	(Sufficient condition for uniform integrability)
	Let $\varphi\ge 0$ be any function such that $\varphi(x) / x \to \infty $ as $x \to  \infty $. If $E \varphi(|X_i|) \le C$ for all $i \in I$, then $\{X_i:  i \in I\} $ is uniformly integrable.
\end{proposition}
\begin{proof}
	\[
		E\left( \frac{|X_i|}{\varphi(|X_i|)} \varphi(|X_i|) ; |X_i| > M \right) 
	.\] 
	As $x \to  \infty $, $\frac{x}{\varphi(x)} \to 0$. So we can apply a supremum bound on it:
	\[
		E\left( \frac{|X_i|}{\varphi(|X_i|)} \varphi(|X_i|) ; |X_i| > M \right) 
		\le \sup \{\frac{x}{\varphi(x)}: x > M\} E\left( \varphi(|X_i| ; |X_i| > M \right)  
		\le \varepsilon_M C
	.\] 
	Since $\varepsilon_M \to  0$ as $M \to \infty $, and the convergence is uniform in $i$, so we have the claim.
\end{proof}



